\section{方法}

\subsection{数据收集}

\subsubsection{数据来源以及数据规模}

研究共采用两个相互独立的数据集对模型进行训练与验证，以充分评估模型架构的可靠性，并实现针对纳米硒的细胞毒性预测。

\begin{enumerate}
  \item 通用数据集：
该数据集来源于 Masarkar 等人 \cite{masarkar2024} 开展的纳米颗粒细胞毒性元分析研究。
数据集汇总了已发表文献中关于金属、金属氧化物、聚合物以及碳基纳米颗粒四大类别的细胞毒性实验数据，是当前纳米材料机器学习毒性预测领域具有代表性的公开数据资源之一。
由于本研究的模型改进工作是在其原有框架基础上展开，因此采用该数据集对基线模型进行复现，以保证模型优化过程具备良好的可比性与有效性。
该数据集原始共包含 4111 条实验记录，经清洗后保留 4102 个有效样本用于后续训练测试划分与模型比较。

  \item 自建数据集：
针对纳米硒这一特定研究对象，本研究从 47 篇已发表文献中系统提取实验数据，构建了纳米硒细胞毒性预测专项数据集。
数据集中共包含 383 个有效样本，每个样本对应 37 个特征参数，涵盖纳米硒的理化性质、实验条件以及生物效应等信息，能够较为全面地反映纳米硒与细胞毒性之间的关联关系。
该数据集为后续改进模型的训练、验证与性能评估提供了核心数据支撑。
\end{enumerate}

\subsubsection{合规与伦理}

本研究所使用的数据均来源于已公开发表的文献和公开数据库，所有数据均可公开获取。
在整理和分析过程中，严格遵守学术规范，对所有来源文献均进行了准确引用，确保数据使用的合法性和透明性。
所有数据均为汇总统计信息，不涉及个人隐私，数据处理过程中遵循科学和负责任的原则。

\subsection{数据预处理}

\subsubsection{数据可行性分析}

本研究自建的纳米硒毒性数据集共包含 383 个样本与 37 个特征。下面从样本规模、数据质量、特征结构及潜在偏差四个维度评估其建模可行性。

1.样本规模与代表性：
数据集覆盖了纳米硒研究中常见的实验条件区间，包括不同的暴露剂量、粒径范围、形态类型与表面修饰策略，能够反映出纳米硒与其理化性质的关系在多个维度上的变化。
383 个样本与 37 个特征的规模虽不及通用纳米毒性数据集，但相对于纳米硒这一较为窄化的研究对象已具备一定代表性，可支持中等复杂度模型的训练，后续将通过正则化、特征筛选与交叉验证等手段控制过拟合风险。

2.数据质量：
所有数据均提取自已发表文献，原始结果可追溯。重复条目、单位不一致、缺失值及明显异常值均已在预处理阶段进行规范化处理。

3.特征结构：
数据集同时包含连续型特征，如粒径、Zeta 电位、暴露浓度，与类别型特征，如细胞系、表面修饰种类、检测方法，构成混合型输入空间。
该结构既为多模态特征融合建模提供了基础，也对类别变量的编码与连续变量的归一化提出了相应要求。

4.潜在限制与风险：
由于数据来源于已有文献，可能存在实验条件偏向或报道选择性偏差，这可能会影响模型在未见实验条件下的泛化能力。
后续模型训练中主要结合交叉验证策略，以降低过拟合和样本偏差带来的影响。

综合以上分析，本数据集在数量、质量、特征维度和建模适用性方面具备较高的可行性，可为纳米硒毒性预测平台的开发提供可靠的数据基础。

\subsubsection{预处理流程}

数据预处理包括以下步骤：

\begin{enumerate}
  \item \textbf{缺失值处理}：由于粒径、Zeta 电位、暴露时间、浓度、细胞存活率诸字段都有若干缺失值，故首先直接剔除含缺失值的记录。此外，“检测方法”字段中若干条目“Characterization”“Western”“Cellular Uptake”属于非毒性检测目的条目，故一并删除。最终共剔除 262 条记录，保留 383 个符合要求的有效样本。

  \item \textbf{多报告值的优先级选择}：参照 Masarkar 等人 \cite{masarkar2024} 提出的方法：同一纳米颗粒在文献中报告多种尺寸时，按无介质中的粒径 > 水相粒径 > 培养基中的粒径的顺序选取，而 Zeta 电位宜取水相测量值，若无该数据则取文献中所给的唯一测量值。

  \item \textbf{单位统一与换算}：把各浓度单位都统一换算为 $\mu$g/mL，凡以 $\mu$M 为单位的浓度都根据硒的分子量（78.96 g/mol）换算为 $\mu$g/mL，同时把暴露时间统一转换为小时（h）。

  \item \textbf{数值特征变换}：由于浓度、时间两变量都呈偏态分布，对二者做对数变换：
        \[
          \text{Log\_Conc} = \ln(1 + \text{Concentration\_μg/mL}), \quad
          \text{Log\_Time} = \ln(1 + \text{Time\_h})
        \]
        同时引入“粒径比”= 水合粒径 / 核心粒径来表征颗粒在水相中的聚集倾向，所有缺失值均按 1.5 填充。

  \item \textbf{分类变量编码}：对表面电荷（正/负/中性）、细胞类型、修饰基团、检测方法、晶型（结晶/无定形）等类别特征，采用标签编码转换为数值，其中表面电荷缺失值填充为“未知”。

  \item \textbf{相关性检查}：对数值特征计算 Pearson 相关系数，未发现 $r > 0.8$ 的高相关度特征对，表明特征冗余度很低。
\end{enumerate}

\subsubsection{数据集特征统计描述}

\paragraph{样本概况}
本研究数据集共包含 383 条有效记录（剔除粒径、Zeta 电位等关键理化性质及细胞存活率缺失记录后），来源于 47 篇已发表文献，时间跨度为 2008--2019 年。其中 2018 年贡献样本量最多（157 条，占 41.0\%），其次为 2019 年（96 条，占 25.1\%）和 2012 年（44 条，占 11.5\%）。所有纳米硒颗粒的形貌均为球形，物种来源以人类细胞为主（371 条，占 96.9\%）。

\paragraph{连续型特征分布}
除表~\ref{tab:statistics} 所列的基本统计量外，我们进一步考察了各特征的分布形态。粒径的分布明显右偏（偏度 3.92），绝大多数颗粒集中在 30 nm 附近，只有极少数样本超过 200 nm，这与纳米硒常规粒径范围十分符合。水合粒径右偏程度较轻（偏度 1.69），其均值（105 nm）明显大于核心粒径均值（65 nm），这表明水化层以及可能存在的弱聚集效应。Zeta 电位分布接近对称（偏度-0.08），范围从 $-60$ mV 到 $+36$ mV，其中负电位占主导，反映出多数表面修饰（如聚糖、PVP）带负电。

浓度分布极端右偏（偏度 8.28），四分位距仅为 $0.52$--$10$ $\mu$g/mL，但最大值高达 $490$ $\mu$g/mL，提示少数实验采用极高剂量，可能对模型产生杠杆效应。暴露时间呈轻微左偏（-0.42），中位数 72 h，说明多数实验采用 24 h 或 72 h 的暴露方案。细胞存活率分布近似均匀（偏度 -0.17），均值与中位数均接近 50\%，意味着数据集正负样本较为平衡，有利于分类模型训练。

为直观展示上述分布，我们绘制了图1（连续型特征的箱线图）和图2（主要类别型特征的频率分布）。

\begin{table}[htbp]
  \centering
  \caption{数据集连续型特征的描述性统计（各特征可用样本量见 $n$ 列）}
  \label{tab:statistics}
  \begin{tabular}{lccccccc}
    \hline
    特征 & $n$ & 均值 & 标准差 & 中位数 & Q1 & Q3 & 偏度 \\
    \hline
    粒径 (nm)            & 339 & 65.30  & 79.72  & 30.00  & 30.00  & 92.15   & 3.92  \\
    水合粒径 (nm)        & 279 & 104.97 & 46.73  & 91.00  & 84.00  & 110.00  & 1.69  \\
    Zeta 电位 (mV)       & 144 & $-5.19$ & 29.76  & $-7.50$ & $-29.10$ & 29.10 & $-0.08$ \\
    暴露时间 (h)         & 363 & 55.21  & 26.14  & 72.00  & 24.00  & 72.00   & $-0.42$ \\
    浓度 ($\mu$g/mL)     & 363 & 12.59  & 38.17  & 0.90   & 0.52   & 10.00   & 8.28  \\
    细胞存活率 (\%)      & 371 & 49.43  & 28.50  & 50.00  & 41.00  & 68.00   & $-0.17$ \\
    \hline
  \end{tabular}
\end{table}

\paragraph{类别型特征分布}
从目前所考察的类别型特征分布来看，数据集涉及 42 种细胞系，结肠癌细胞 HCT 116 出现频率最高（98 次，25.6\%），其次为肝癌细胞 HepG2（41 次，10.7\%）、黑色素瘤细胞 A375（23 次，6.0\%）和乳腺癌细胞 MCF-7（23 次，6.0\%），前 10 种细胞系合计覆盖 71.0\% 的样本。组织来源以结肠（28.7\%）、肝脏（16.7\%）和脑组织（8.9\%）为主。

表面修饰类型共 43 种，其中聚糖修饰（PTR）占比最高（127 次，33.2\%），其后依次为聚糖蛋白复合物 PR（18 次，4.7\%）和 PAMAM+RGD+siRNA 复合修饰（15 次，3.9\%）。合成方法以化学还原法为主（352 次，占 91.9\%），毒性检测方法中最常见的是 MTT 法（206 次，53.8\%）。

\subsubsection{数据局限性}
本研究存在以下局限性：（1）数据集规模有限（383 条有效记录），来源于文献手工提取，约 40.6\% 的原始记录因粒径、Zeta 电位等关键理化性质及细胞存活率字段缺失而被剔除，且可能存在发表偏倚，即阴性结果或无显著毒性的数据较少被报道；（2）不同文献的实验条件（如毒性检测方法涉及 MTT、MTS、CCK-8 等不同试剂）存在异质性，可能引入系统性噪声；（3）数据集时间跨度为 2008--2019 年，早期文献的实验方法和报告标准与近期存在差异。未来研究可通过多中心数据共享、标准化实验方案以及主动学习策略来扩充数据规模并降低异质性。

\subsubsection{特征设置与筛选}

基础特征包含以下几类信息：暴露条件（浓度、时间）、颗粒理化属性（粒径、Zeta 电位、尺寸比、表面电荷）、生物体系属性（细胞类型）以及实验条件（检测方法、修饰基团）。针对硒纳米颗粒专用建模，在基础特征之外进一步引入 5 项增强特征，即 IC50、元素 Se 含量、长期稳定性、结合能 Se3d 和晶型编码。

为评估新增材料化学属性对模型性能的提升作用，保持基础特征不变，对上述 5 项增强特征的全部非空组合进行穷举搜索，共 31 种组合，并以三折交叉验证 MAE 作为筛选指标，用于确定最优增强特征子集。

\subsubsection{模型架构}
本研究围绕表格型纳米毒性数据比较了树模型、注意力增强树模型与深度表格模型三类方法。在基准数据集上，纳入随机森林、XGBoost、ABRF、LARF、TabTransformer 和蒸馏 TabTransformer 六种模型，以系统评估不同架构对异质纳米颗粒数据的适应性；在硒纳米颗粒专用数据集上，重点比较基础特征与增强特征条件下各模型的性能差异。

随机森林作为核心基线模型，用于评估中小样本场景下的稳健性与可解释性；XGBoost 用于测试梯度提升树在增强特征条件下的高精度建模能力；ABRF 与 LARF 用于考察注意力机制引入后是否能够提升树模型表现；TabTransformer 与蒸馏 TabTransformer 则作为深度表格模型代表，用于与传统树模型进行对照。

\subsubsection{训练与验证策略}
对于大规模多类型纳米颗粒基准数据集，采用 7:3 训练测试划分，以测试集性能作为主要比较依据。回归任务使用 $R^2$、MAE 和 RMSE 进行评价，分类任务使用 Accuracy、F1 和 AUC 进行评价。

对于硒纳米颗粒专用数据集，采用三折交叉验证，以 CV MAE 和全量 $R^2$ 作为主要回归评价指标，并用于比较基础特征与增强特征两种设置下不同模型的表现。增强特征筛选阶段同样采用三折交叉验证 MAE 作为统一指标，以保证不同特征组合之间的结果可比性。